% !TEX root = ../../ctfp-print.tex

\lettrine[lhang=0.17]{我}{意识到, 我们或许正在远离编程}, 一头扎进硬核数学里去. 可你永远说不
准下一场编程大革命会带来什么, 也说不准需要什么样的数学才能理解它. 眼下正流传着一些非常有意思的想法,
比如带连续时间的函数式反应式编程, 用依赖类型来扩展 Haskell 的类型系统, 或者把同伦类型论探索进编程里.

到目前为止, 我一直随随便便地把类型等同于取值的\emph{集合}. 严格说来这并不正确, 因为这种处理方式没有
考虑到一个事实: 在编程里, 我们是\emph{计算}出取值的, 而计算是一个要花时间的过程, 在
极端情况下还可能不终止. 发散的计算是每一种图灵完备语言的一部分.

还有一些根本性的理由, 说明集合论未必是充当计算机科学 —— 甚至数学自身 —— 之基础的最佳选择. 一个好的
类比是: 集合论就像与某种特定体系结构绑定的汇编语言. 如果你想在不同的体系结构上跑你的数学, 就得使
用更通用的工具.

一种可能是用空间来代替集合. 空间带有更多结构, 而且可以不诉诸集合来定义. 一说到空间, 人们通常就
会联想到拓扑, 而连续性这类东西正是靠拓扑才能定义的. 至于处理拓扑的常规途径, 你猜对了, 正是通过集合论.
尤其地, 子集这个概念在拓扑里居于核心地位. 不出意料, 范畴论者把这个想法推广到了 $\Set$ 之外的范畴.
有一类范畴, 恰好具备充当集合论替代品所需的一切性质, 它叫作 \newterm{topos 拓扑斯} (复数: topoi);
除别的东西之外, 它还提供了一个推广了的子集概念.

\section{子物件分类子}

我们先试着不用元素、而用函数来表达子集这个概念. 从某个集合 $a$ 到 $b$ 的任何一个函数 $f$ 都定义了
$b$的一个子集 —— 也就是 $a$ 在 $f$ 下的像. 但定义同一个子集的函数有很多. 我们需要更精确一些. 首先,
我们可以把注意力放在单射的函数上 —— 也就是不会把多个元素压成一个的那种. 单射函数把一个集合
``单射地''送入另一个集合. 对有限集合, 你不妨把单射函数想象成一组平行箭头, 把一个集合的元素连到
另一个集合的元素上. 当然, 第一个集合不可能比第二个集合更大, 否则那些箭头就必然要汇聚. 可这里
还残留着一些歧义: 可能还有另一个集合 $a'$, 以及另一个从它到 $b$ 的单射函数 $f'$, 挑中的是同一个子集.
但你很容易说服自己, 这样的集合必定与 $a$ 同构. 我们可以用这个事实把子集定义成一族单射函数,
它们之间由各自定义域之间的同构相联系. 更精确地说, 我们说下面两个单射函数
\begin{align*}
  f  & \Colon a \to b  \\
  f' & \Colon a' \to b
\end{align*}
是等价的, 若存在一个同构
\[h \Colon a \to a'\]
使得:
\[f = f'\ .\ h\]
这样一族等价的单射定义出 $b$ 的一个子集.

\begin{figure}[H]
  \centering
  \includegraphics[width=0.4\textwidth]{images/subsetinjection.jpg}
\end{figure}

\noindent
只要我们把这些单射函数换成单态射, 这个定义就能抬升到任意范畴. 提醒一下, 从 $a$ 到 $b$ 的单态射 $m$
是用它的普遍性质定义的. 对任何物件 $c$ 和任何一对态射
\begin{align*}
  g  & \Colon c \to a \\
  g' & \Colon c \to a
\end{align*}
只要
\[m\ .\ g = m\ .\ g'\]
那就必定有 $g = g'$.

\begin{figure}[H]
  \centering
  \includegraphics[width=0.4\textwidth]{images/monomorphism.jpg}
\end{figure}

\noindent
在集合上, 如果我们想一想一个函数 $m$ \emph{不}是单态射意味着什么, 这个定义就更容易理解了. 它会把
$a$的两个不同元素映到 $b$ 的同一个元素上. 于是我们能找到两个函数 $g$ 与 $g'$, 它们只在
那两个元素上有区别. 再与 $m$ 做后复合, 这个区别就会被掩盖掉.

\begin{figure}[H]
  \centering
  \includegraphics[width=0.4\textwidth]{images/notmono.jpg}
\end{figure}

\noindent
定义子集还有另一种办法: 用一个函数, 叫作特征函数. 它是一个从集合 $b$ 到两元素集合 $\Omega$ 的
函数$\chi$. 这个集合里的一个元素被指定为 ``真'', 另一个被指定为 ``假''. 这个函数给 $b$ 里
那些属于该子集的元素指派 ``真'', 给不属于的指派 ``假''.

剩下要说明的, 是把 $\Omega$ 的一个元素指定为 ``真'' 究竟是什么意思. 我们可以用标准的那套把戏:
用一个从单元素集合到 $\Omega$ 的函数. 我们把这个函数叫作 $true$:
\[true \Colon 1 \to \Omega\]

\begin{figure}[H]
  \centering
  \includegraphics[width=0.4\textwidth]{images/true.jpg}
\end{figure}

\noindent
这些定义可以这样组合起来: 让它们不仅定义出什么是子物件, 还能在不谈论元素的情况下定义出那个特殊的
物件$\Omega$. 思路是, 我们希望态射 $true$ 代表一个 ``一般的'' 子物件. 在 $\Set$ 里, 它从
两元素集合$\Omega$ 中挑出一个单元素子集. 这已经一般得不能再一般了. 它显然是一个真子集, 因为
$\Omega$还有一个元素 \emph{不}在那个子集里.

在更一般的设定下, 我们把 $true$ 定义成一个从终止物件到 \emph{分类物件} $\Omega$ 的单态射.
但我们得先定义分类物件. 我们需要一条普遍性质, 把这个物件与特征函数联系起来. 结果发现, 在 $\Set$ 里,
沿特征函数 $\chi$ 对 $true$ 取拉回, 既定义出子集 $a$, 又定义出把它嵌入 $b$ 的那个单射函数. 下面就是
那个拉回图:

\begin{figure}[H]
  \centering
  \includegraphics[width=0.4\textwidth]{images/pullback.jpg}
\end{figure}

\noindent
我们来分析这张图. 拉回方程是:
\[true\ .\ unit = \chi\ .\ f\]
函数 $true\ .\ unit$ 把 $a$ 的每个元素都映到 ``真''. 因此 $f$ 必定把 $a$ 的所有元素都映到 $b$ 中
那些使 $\chi$ 为 ``真'' 的元素上. 按定义, 这些元素就是特征函数 $\chi$ 所指定的那个子集的元素. 所以
$f$ 的像确实就是我们所谈的那个子集. 拉回的普遍性则保证 $f$ 是单射的.

这张拉回图可以用来在 $\Set$ 之外的范畴里定义分类物件. 这样的范畴必须有一个终止物件, 好让
我们定义单态射 $true$. 它还必须拥有拉回 —— 实际的要求是它必须拥有所有有限极限 (拉回就是有限极限的
一个例子). 在这些假设下, 我们用如下性质来定义分类物件 $\Omega$: 对每个单态射 $f$, 都存在唯一的
态射$\chi$, 把这张拉回图补全.

我们来分析最后这句话. 构造一个拉回时, 我们拿到的是三个物件 $\Omega$, $b$ 和 $1$, 以及两个态射
$true$与 $\chi$. 拉回的存在意味着, 我们能找到最好的那个物件 $a$, 配上两个态射 $f$ 与 $unit$ (后者由
终止物件的定义唯一确定), 使得图交换.

而在这里, 我们解的是另一套方程. 我们在变动 $a$ \emph{和} $b$ 的同时, 去解出 $\Omega$ 与 $true$.
给定$a$ 与 $b$, 可能有一个单态射 $f \Colon a \to b$, 也可能没有. 但若有, 我们希望它是某个 $\chi$ 的
拉回. 而且我们还希望这个 $\chi$ 由 $f$ 唯一确定.

我们不能说单态射 $f$ 与特征函数 $\chi$ 之间存在一一对应, 因为拉回只在同构的意义下唯一.
但请回想我们先前把子集定义成一族等价的单射. 我们可以把它推广: 把 $b$ 的一个子物件定义成一族到 $b$
的等价单态射. 这族单态射与我们这张图的那族等价拉回一一对应.

于是我们能把 $b$ 的子物件构成的集合 $Sub(b)$ 定义成一族单态射, 并看出它同构于从 $b$ 到 $\Omega$ 的
态射集合:
\[Sub(b) \cong \cat{C}(b, \Omega)\]
这恰好是两个函子之间的一个自然同构. 换句话说, $Sub(-)$ 是一个可表的 (反变) 函子, 它的表示就是
物件$\Omega$.

\section{拓扑斯}

一个拓扑斯是这样的范畴:

\begin{enumerate}
  \tightlist
  \item
是笛卡尔闭的: 它拥有所有的积、终止物件, 以及指数 (指数定义成积的右伴随),
  \item
拥有所有有限图式的极限,
  \item
拥有一个子物件分类子 $\Omega$.
\end{enumerate}

这套性质让拓扑斯在大多数应用里都能当仁不让地顶替 $\Set$. 它还有一些由定义推出来的额外性质. 比如,
一个拓扑斯拥有所有有限余极限, 包括初始物件.

把子物件分类子定义成终止物件的两个副本的余积 (和) 会很诱人 —— 在 $\Set$ 里它正是如此 ——
但我们想要更一般的情形. 满足这一点的拓扑斯叫作布尔的.

\section{拓扑斯与逻辑}

在集合论里, 一个特征函数可以被解读为定义集合元素的一条性质 —— 即一个对某些元素为真、对另一些元素为
假的 \newterm{predicate 谓词}. 谓词 $isEven$ 从自然数集合里挑出偶数子集. 在拓扑斯里, 我们可以把
谓词的概念推广成从物件 $a$ 到 $\Omega$ 的一个态射. 这就是 $\Omega$ 有时被叫作真值物件的原因.

谓词是逻辑的积木. 一个拓扑斯含有研究逻辑所需的全部器械. 它的积对应于逻辑合取 (逻辑 \emph{与}),
余积对应于析取 (逻辑 \emph{或}), 指数则对应于蕴涵. 逻辑的所有标准公理在拓扑斯里都成立, 唯独排中律
(或者等价的双重否定消去) 例外. 这就是为什么拓扑斯的逻辑对应于构造性逻辑或直觉主义逻辑.

直觉主义逻辑一直在稳步扩大地盘, 并从计算机科学那里获得了意想不到的支持. 经典的排中律观念建立在
一种信念上: 存在绝对的真理 —— 任何陈述要么为真, 要么为假; 或者像古罗马人所说,
\emph{tertium non datur} (没有第三种选择). 但我们能知道一件事为真还是为假的唯一途径, 就是
我们能证明它或反驳它. 证明是一个过程、一次计算 —— 而我们知道, 计算是要花时间和资源的. 有些时候,
计算可能永远不终止. 如果我们无法在有限时间内证明一个陈述, 那么声称它为真就毫无意义. 拓扑斯有
着更细致的真值物件, 从而为给有趣的逻辑建模提供了一个更一般的框架.

\section{挑战}

\begin{enumerate}
  \tightlist
  \item
证明: 沿特征函数对 $true$ 取拉回所得到的那个函数 $f$ 必定是单射的.
\end{enumerate}
